% =============================================================================
% Chapter 3 — System Description and Modeling
% Thesis: Modern Control Methods for Agricultural Irrigation
% Author: Tara Torbati, ITMO University
% =============================================================================

\section{System Description and Modeling}
\label{ch:system}

This chapter presents the complete mathematical and physical description of the agricultural system studied in this thesis. The objective is to construct a high-fidelity simulation environment that serves a dual purpose: first, as the virtual plant on which the Model Predictive Controller operates, and second, as the training environment for the Reinforcement Learning agent developed in Chapter~\ref{ch:controllers}. The modeling framework builds upon the agent-based model (ABM) proposed by Lopez-Jimenez et al.~\cite{Lopezjimenez2024}, which couples soil water dynamics with crop growth on a topographically heterogeneous field. However, several critical modifications and extensions are introduced to adapt the model to the specific agronomic conditions of rice cultivation in Gilan Province, Iran, and to correct certain mass-balance inconsistencies identified in the original formulation.

The chapter is organized as follows. Section~\ref{sec:field} describes the physical field site and the digital elevation model used to define the agent grid. Section~\ref{sec:abm_equations} presents the complete ABM equations. Section~\ref{sec:mass_correction} details the mass conservation correction applied to the original water exchange formulation. Section~\ref{sec:climate} describes the climate data sources and the two contrasting scenarios used throughout this thesis. Section~\ref{sec:soil_crop} specifies the soil and crop parameters. Section~\ref{sec:water_budget} derives the agronomically grounded water budget. Finally, Section~\ref{sec:validation} presents the ABM validation results, including cross-validation against satellite soil wetness observations.


% =============================================================================
\section{Field Description and Digital Elevation Model}
\label{sec:field}
% =============================================================================

The study site is located in Gilan Province in northwestern Iran, at coordinates 38.298°N, 48.847°E. Gilan is one of Iran's principal rice-producing regions, situated along the southwestern coast of the Caspian Sea. The province experiences a humid subtropical climate with hot summers and relatively mild winters, though seasonal precipitation can vary dramatically between years. Despite its comparatively favorable water endowment relative to Iran's central plateau, Gilan faces mounting pressures from declining aquifer levels and increasing competition between agricultural, urban, and industrial water users~\cite{nouri2023,mesgaran2018}. These conditions make it a representative case study for evaluating constrained irrigation control strategies.

The selected field covers approximately 6 hectares and is discretized into a $10 \times 13$ rectangular grid, yielding $N = 130$ crop--soil agents. This resolution represents a practical balance between spatial detail and computational tractability: each agent corresponds to a plot of roughly 460~m$^2$, which is sufficiently fine to capture the field's topographical heterogeneity while remaining computationally feasible for the MPC optimization solved at every daily time step.

The topographical data were obtained from the USGS Earth Explorer digital elevation database, which provides Shuttle Radar Topography Mission (SRTM) data at 30-meter resolution. The raw elevation raster was processed in QGIS~3.34 to extract the study region, reproject it to UTM Zone~38N, and resample it to the $10 \times 13$ agent grid. The resulting digital elevation model (DEM) is stored as a GeoTIFF file and loaded programmatically by the simulation environment.

\begin{figure}[htbp]
    \centering
    \includegraphics[width=0.75\textwidth]{figures/Figure_1.png}
    \caption{Three-dimensional surface plot of the 130-agent grid. Elevation ranges from 70~m to 181~m above sea level, forming a bowl-shaped topology with water convergence toward the southeastern corner.}
    \label{fig:topo3d}
\end{figure}

The elevation field exhibits a distinctive bowl-shaped topology, with elevations ranging from 70~m in the southeastern corner to 181~m in the northwestern and southwestern ridges. This 111~m elevation gradient creates significant lateral water redistribution through surface runoff: agents on higher ground shed excess water to their downhill neighbors, while agents in the valley bottom accumulate inflows from multiple directions. This topographical heterogeneity is a defining feature of the study site and the principal motivation for adopting an agent-based rather than a lumped-parameter modeling approach.


\subsection{Directed Graph Construction}
\label{sec:directed_graph}

The topographical information is encoded in a directed graph that governs lateral water exchange between agents. Each agent occupies a cell in the $10 \times 13$ grid and is assigned a normalized elevation $\gamma^{(n)} \in [0, 1]$, computed by min--max normalization of the raw DEM values:
\begin{equation}
    \gamma^{(n)} = \frac{z^{(n)} - z_{\min}}{z_{\max} - z_{\min}}
    \label{eq:gamma}
\end{equation}
where $z^{(n)}$ is the raw elevation of agent $n$. An agent at position $(r, c)$ can exchange water with up to eight neighbors in its Moore neighborhood (the eight surrounding cells). A directed edge from agent $n$ to agent $m$ is created if and only if $\gamma^{(m)} < \gamma^{(n)}$, so that water flows strictly downhill. The number of lower-elevation neighbors for each agent is denoted $N_r^{(n)}$.

\begin{figure}[htbp]
    \centering
    \includegraphics[width=0.75\textwidth]{figures/Figure_2.png}
    \caption{Water flow directions between agents, overlaid on the normalized elevation map. Arrows indicate the direction of lateral surface runoff, governed by the directed graph.}
    \label{fig:flow_arrows}
\end{figure}

Figure~\ref{fig:flow_arrows} visualizes the resulting directed graph. The arrow density reflects the terrain gradient: steep areas near the ridges exhibit dense, nearly parallel flow lines, while the flat valley floor shows sparser, more diffuse connections. Agents in the southeastern corner with $N_r = 0$ function as hydrological sinks---they have no lower neighbors and therefore retain all accumulated water. These sink agents are particularly susceptible to waterlogging stress, a phenomenon that the ABM captures through the waterlogging stress function $h_6$ described in Section~\ref{sec:stress_growth}.

\begin{figure}[htbp]
    \centering
    \includegraphics[width=0.7\textwidth]{figures/Figure_3.png}
    \caption{Spatial distribution of the neighbor count $N_r$ for each agent. Red-bordered cells indicate sink agents ($N_r = 0$) that accumulate runoff from uphill neighbors.}
    \label{fig:sink_agents}
\end{figure}

Three sink agents were identified, all located in the southeastern corner of the field at the lowest elevations (Figure~\ref{fig:sink_agents}). This spatial concentration of sinks creates a natural drainage basin structure that any irrigation controller must account for: over-irrigating uphill agents will cause water to pool at these sinks, simultaneously wasting budget and inducing waterlogging stress in the valley.


% =============================================================================
\section{Agent-Based Model Equations}
\label{sec:abm_equations}
% =============================================================================

The crop--soil system is modeled as a set of $N = 130$ coupled discrete-time dynamical systems, one per agent, following the ABM framework introduced by Lopez-Jimenez et al.~\cite{Lopezjimenez2024}. Each agent $n \in \{1, \ldots, N\}$ maintains four state variables that evolve at a daily time step $k$. These states capture the essential dynamics of the soil--plant--atmosphere continuum: water availability, thermal accumulation, crop maturation, and biomass production.

The four state equations governing the evolution of agent $n$ are:
\begin{align}
    x_{1,k+1}^{(n)} &= x_{1,k}^{(n)} - \phi_{1,k}^{(n)} - \phi_{2,k}^{(n)} - \phi_{3,k}^{(n)} - \phi_{4,\text{out},k}^{(n)} + \phi_{4,\text{in},k}^{(n)} + u_{e,k}^{(1)} + u_k^{(n)} \label{eq:x1} \\
    x_{2,k+1}^{(n)} &= x_{2,k}^{(n)} + h_{1,k}^{(n)} \label{eq:x2} \\
    x_{3,k+1}^{(n)} &= x_{3,k}^{(n)} + \theta_{11}\bigl(1 - h_{2,k}^{(n)}\bigr) + \theta_{12}\bigl(1 - h_{3,k}^{(n)}\bigr) \label{eq:x3} \\
    x_{4,k+1}^{(n)} &= x_{4,k}^{(n)} + \theta_{13} \cdot h_{3,k}^{(n)} \cdot h_{6,k}^{(n)} \cdot h_{7,k}^{(n)} \cdot g_k^{(n)} \cdot u_{e,k}^{(4)} \label{eq:x4}
\end{align}
where $x_1$ is the volumetric soil water content in the root zone (mm), $x_2$ is the cumulative thermal time (°C$\cdot$day), $x_3$ is the cumulative maturation index incorporating heat and drought stress effects, and $x_4$ is the aboveground dry biomass (g/m$^2$). The environmental inputs are rainfall $u_e^{(1)}$ (mm/day), reference evapotranspiration $u_e^{(2)}$ (mm/day), mean air temperature $u_e^{(3)}$ (°C), solar radiation $u_e^{(4)}$ (MJ/m$^2$/day), and maximum daily temperature $u_e^{(6)}$ (°C). The control input $u^{(n)}_k$ is the irrigation applied to agent $n$ on day $k$ (mm/day). Each flux and stress function is defined below.

Note that Equation~\eqref{eq:x1} includes two water exchange terms, $\phi_{4,\text{in}}$ and $\phi_{4,\text{out}}$, rather than the single $\phi_4$ term in the original formulation by Lopez-Jimenez et al.~\cite{Lopezjimenez2024}. This modification corrects a mass conservation error discussed in detail in Section~\ref{sec:mass_correction}. Additionally, Equation~\eqref{eq:x4} includes the drought stress multiplier $h_3$, which was not present in the original biomass equation; this correction is also described in Section~\ref{sec:mass_correction}.

Table~\ref{tab:notation} summarizes the key variables and parameters used throughout this chapter.

\begin{table}[htbp]
\centering
\caption{General notation for the ABM state variables, fluxes, and inputs.}
\label{tab:notation}
\begin{tabular}{cl}
\toprule
\textbf{Symbol} & \textbf{Description} \\
\midrule
$N$ & Number of agents (130) \\
$k$ & Daily time step index \\
$x_1^{(n)}$ & Soil water content in the root zone (mm) \\
$x_2^{(n)}$ & Cumulative thermal time (°C$\cdot$day) \\
$x_3^{(n)}$ & Cumulative maturation stress index \\
$x_4^{(n)}$ & Aboveground dry biomass (g/m$^2$) \\
$u^{(n)}$ & Irrigation input for agent $n$ (mm/day) \\
$u_e^{(1)}$ & Rainfall (mm/day) \\
$u_e^{(2)}$ & Reference evapotranspiration ET$_0$ (mm/day) \\
$u_e^{(3)}$ & Mean air temperature (°C) \\
$u_e^{(4)}$ & Solar radiation (MJ/m$^2$/day) \\
$u_e^{(6)}$ & Maximum daily temperature (°C) \\
$\phi_1$ & Crop transpiration (mm/day) \\
$\phi_2$ & Surface runoff (mm/day) \\
$\phi_3$ & Deep drainage (mm/day) \\
$\phi_{4,\text{in}}, \phi_{4,\text{out}}$ & Lateral water exchange: inflow, outflow (mm/day) \\
$\gamma^{(n)}$ & Normalized elevation of agent $n$ \\
$N_r^{(n)}$ & Number of lower-elevation neighbors \\
\bottomrule
\end{tabular}
\end{table}


\subsection{Soil Water Fluxes}
\label{sec:soil_fluxes}

\paragraph{Crop transpiration} ($\phi_1$) represents the water extracted from the soil by plant roots to meet evaporative demand. It is modeled as a demand-limited process:
\begin{equation}
    \phi_{1,k}^{(n)} = \min\!\Big(\theta_1\bigl(x_{1,k}^{(n)} - \theta_2 \theta_5\bigr),\; u_{e,k}^{(2)}\Big)
    \label{eq:phi1}
\end{equation}
where $\theta_1$ is the water uptake coefficient, $\theta_2$ is the volumetric wilting point, and $\theta_5$ is the effective root zone depth (mm). The transpiration cannot exceed the atmospheric demand given by the reference evapotranspiration ET$_0$, and it drops to zero when soil water reaches the wilting point.

\paragraph{Surface runoff} ($\phi_2$) is computed using the SCS curve number method, which partitions rainfall between infiltration and overland flow:
\begin{equation}
    \phi_{2,k} = 
    \begin{cases}
        \dfrac{\bigl(u_{e,k}^{(1)} - \theta_3\bigr)^2}{u_{e,k}^{(1)} + 4\,\theta_3}, & u_{e,k}^{(1)} > \theta_3 \\[8pt]
        0, & u_{e,k}^{(1)} \leq \theta_3
    \end{cases}
    \label{eq:phi2}
\end{equation}
where $\theta_3$ is the initial abstraction threshold (mm). Runoff is triggered only when rainfall exceeds this threshold, and the nonlinear response saturates for heavy rainfall events.

\paragraph{Deep drainage} ($\phi_3$) represents water percolating below the root zone and is activated when soil water exceeds field capacity:
\begin{equation}
    \phi_{3,k}^{(n)} = 
    \begin{cases}
        \theta_4\bigl(x_{1,k}^{(n)} + u_{e,k}^{(1)} - \phi_{2,k} - \theta_6 \theta_5\bigr), & \text{if } x_{1,k}^{(n)} + u_{e,k}^{(1)} - \phi_{2,k} > \theta_6 \theta_5 \\[4pt]
        0, & \text{otherwise}
    \end{cases}
    \label{eq:phi3}
\end{equation}
where $\theta_4$ is the drainage coefficient governing the rate of percolation, and $\theta_6 \theta_5$ is the total field capacity of the root zone (mm). In the silty loam soils of this study, $\theta_4 = 0.05$, indicating relatively slow drainage consistent with the moderate hydraulic conductivity of this soil type.

\paragraph{Lateral water exchange} ($\phi_{4,\text{in}}$ and $\phi_{4,\text{out}}$) represents the topography-driven redistribution of excess soil water between adjacent agents. The formulation adopted in this thesis corrects the original model and is detailed in Section~\ref{sec:mass_correction}.


\subsection{Stress Functions and Crop Growth}
\label{sec:stress_growth}

\paragraph{Thermal time accumulation} ($h_1$) drives phenological development. Only temperatures above the crop's base temperature $\theta_7$ contribute to development:
\begin{equation}
    h_{1,k} = \max\!\bigl(u_{e,k}^{(3)} - \theta_7,\; 0\bigr)
    \label{eq:h1}
\end{equation}
For rice, $\theta_7 = 10$°C, reflecting the crop's tropical origin and its requirement for warm conditions to initiate growth. The thermal time state $x_2$ accumulates approximately 10--15~°C$\cdot$day during the summer months in Gilan, reaching roughly 1200--1500~°C$\cdot$day over the 120-day season.

\paragraph{Heat stress} ($h_2$) models the reduction in radiation use efficiency when maximum daily temperature exceeds a crop-specific threshold:
\begin{equation}
    h_{2,k} = 
    \begin{cases}
        1, & u_{e,k}^{(6)} \leq \theta_9 \\[4pt]
        1 - \dfrac{u_{e,k}^{(6)} - \theta_9}{\theta_{10} - \theta_9}, & \theta_9 < u_{e,k}^{(6)} \leq \theta_{10} \\[8pt]
        0, & u_{e,k}^{(6)} > \theta_{10}
    \end{cases}
    \label{eq:h2}
\end{equation}
For rice, $\theta_9 = 35$°C and $\theta_{10} = 42$°C. Maximum daily temperatures in Gilan during the April--August rice season rarely exceed 35°C, so heat stress plays a minor role in this study.

\paragraph{Drought stress} ($h_3$) quantifies the reduction in crop performance when transpiration falls short of atmospheric demand:
\begin{equation}
    h_{3,k}^{(n)} = 1 - \theta_{14} \cdot h_{4,k}^{(n)}
    \label{eq:h3}
\end{equation}
where
\begin{equation}
    h_{4,k}^{(n)} = 
    \begin{cases}
        1 - \dfrac{\phi_{1,k}^{(n)}}{u_{e,k}^{(2)}}, & \phi_{1,k}^{(n)} < u_{e,k}^{(2)} \\[8pt]
        0, & \phi_{1,k}^{(n)} \geq u_{e,k}^{(2)}
    \end{cases}
    \label{eq:h4}
\end{equation}
The parameter $\theta_{14}$ controls the crop's sensitivity to drought. For rice, $\theta_{14} = 0.8$, indicating very high sensitivity---when transpiration drops to zero, $h_3$ falls to 0.2, reducing biomass accumulation by 80\%. This extreme sensitivity reflects rice's biological requirement for continuous water availability and is the fundamental reason why rice was selected as the primary crop for this study: it provides the most demanding test case for any irrigation controller.

\paragraph{Waterlogging stress} ($h_6$) penalizes biomass growth when soil water exceeds field capacity, simulating oxygen deprivation in the root zone:
\begin{equation}
    h_{6,k}^{(n)} = \text{clip}\!\left(1 - \frac{x_{1,k}^{(n)} - \text{FC}}{\text{FC}},\; 0,\; 1\right)
    \label{eq:h6}
\end{equation}
where $\text{FC} = \theta_6 \theta_5$ is the total field capacity. This stress function is particularly relevant for the sink agents in the southeastern corner of the field, which accumulate lateral inflow from multiple uphill neighbors.

\paragraph{Low temperature stress} ($h_7$) prevents biomass accumulation when mean temperature falls below the base temperature:
\begin{equation}
    h_{7,k} = 
    \begin{cases}
        1, & u_{e,k}^{(3)} > \theta_7 \\
        0, & u_{e,k}^{(3)} \leq \theta_7
    \end{cases}
    \label{eq:h7}
\end{equation}

\paragraph{Growth function} ($g$) determines the fraction of intercepted radiation that contributes to biomass, following a sigmoid model driven by cumulative thermal time:
\begin{equation}
    g_k^{(n)} = 
    \begin{cases}
        \dfrac{\theta_{19}}{1 + \exp\!\bigl(-0.01(x_{2,k}^{(n)} - \theta_{20})\bigr)}, & x_{2,k}^{(n)} \leq \dfrac{\theta_{18}}{2} \\[10pt]
        \dfrac{\theta_{19}}{1 + \exp\!\bigl(0.01(x_{2,k}^{(n)} + x_{3,k}^{(n)} - \theta_{18})\bigr)}, & x_{2,k}^{(n)} > \dfrac{\theta_{18}}{2}
    \end{cases}
    \label{eq:growth}
\end{equation}
The first branch governs the vegetative phase, during which radiation interception increases sigmoidally as leaf area expands. The second branch governs senescence, during which interception declines as the accumulated stress state $x_3$ pushes the crop toward maturity. The parameter $\theta_{19} = 0.95$ represents the maximum fraction of radiation that the canopy can intercept, $\theta_{18} = 1800$~°C$\cdot$day is the thermal time to full maturity, and $\theta_{20} = 600$~°C$\cdot$day is the thermal time at which radiation interception reaches 50\%.

It is critical to note that the first branch uses $x_2$ (thermal time) rather than $x_3$ (maturation stress) as the sigmoid driver: $x_2$ accumulates on the order of 10~°C/day for rice, reaching approximately 1200 over 120 days, while $x_3$ accumulates only about 0.7 over the same period and would effectively freeze the sigmoid at its minimum.


% =============================================================================
\section{Mass Conservation Corrections}
\label{sec:mass_correction}
% =============================================================================

During the implementation and validation of the ABM, two significant modeling deficiencies were identified in the original formulation of Lopez-Jimenez et al.~\cite{Lopezjimenez2024}. Both were corrected in the present implementation. These corrections constitute an original contribution of this thesis.


\subsection{Water Exchange Mass Balance}
\label{sec:phi4_correction}

In the original model, the lateral water exchange function $\phi_4$ appears only as an inflow term in the soil water mass balance (Equation~1a of~\cite{Lopezjimenez2024}). When an agent's soil water exceeds field capacity, the excess is divided equally among its lower-elevation neighbors and added to their soil water. However, the corresponding outflow is \emph{not} subtracted from the sending agent's state. This omission violates mass conservation: water is created at every exchange event, as the sender retains its original water content while simultaneously transferring a copy to each receiver.

The corrected formulation separates the exchange into explicit inflow and outflow components. For each agent $n$ with soil water exceeding field capacity, the total outflow is:
\begin{equation}
    \phi_{4,\text{out},k}^{(n)} = 
    \begin{cases}
        x_{1,k}^{(n)} - \theta_6 \theta_5, & x_{1,k}^{(n)} > \theta_6 \theta_5 \\
        0, & \text{otherwise}
    \end{cases}
    \label{eq:phi4out}
\end{equation}
This outflow is divided equally among the $N_r^{(n)}$ lower-elevation neighbors, so each receiver $m \in \texttt{sends\_to}(n)$ gains:
\begin{equation}
    \phi_{4,\text{in},k}^{(m)} \mathrel{+}= \frac{\phi_{4,\text{out},k}^{(n)}}{N_r^{(n)}}
    \label{eq:phi4in}
\end{equation}
Critically, the full outflow $\phi_{4,\text{out}}$ is now subtracted from the sender's mass balance in Equation~\eqref{eq:x1}. This ensures that water is conserved across the entire field: the total water entering the system (rainfall plus irrigation) minus the total water leaving (transpiration, deep drainage, and surface runoff at the field boundary) exactly equals the change in total soil water storage.

The practical impact of this correction is significant. Without it, the ABM accumulates phantom water over the 120-day simulation, producing unrealistically high soil moisture levels and masking the effects of drought stress. With the correction, the soil water dynamics exhibit the expected behavior: gradual depletion between rainfall events, responsive recovery following irrigation, and spatial differentiation between hilltop and valley agents.


\subsection{Drought Stress in the Biomass Equation}
\label{sec:h3_correction}

In the original model, the drought stress function $h_3$ appears only in the maturation equation ($x_3$, Equation~1c of~\cite{Lopezjimenez2024}), where it accelerates crop senescence under water deficit. However, $h_3$ is absent from the biomass equation ($x_4$, Equation~1d). This means that a completely unirrigated crop would accumulate the same daily biomass increment as a fully irrigated one, provided all other conditions were equal. The only effect of drought would be to advance the onset of senescence, which reduces biomass accumulation only indirectly and with considerable delay.

This omission was identified during the initial ABM validation, when the no-irrigation baseline produced biomass values that were nearly indistinguishable from irrigated scenarios for wheat. The corrected formulation multiplies the radiation use efficiency term in Equation~\eqref{eq:x4} by $h_3$, so that drought directly reduces the daily biomass increment. This modification follows the approach of the FAO AquaCrop model, in which water stress reduces radiation use efficiency proportionally to the ratio of actual to potential transpiration. With $\theta_{14} = 0.8$ for rice, maximum drought reduces the biomass increment by up to 80\%, making irrigation decisions directly consequential for yield outcomes.

This correction is particularly important for the controller comparison objective of this thesis. Without it, the ABM would produce only marginal differences between irrigated and unirrigated outcomes, rendering any comparison between MPC, RL, and baseline controllers meaningless. The original crop used in the baseline paper was wheat in Colombia, which is less water-sensitive (lower $\theta_{14}$); it is possible that the omission went undetected in that context because the irrigation response was already small.


% =============================================================================
\section{Climate Data and Scenarios}
\label{sec:climate}
% =============================================================================

Daily meteorological data for the study site were obtained from the NASA POWER (Prediction of Worldwide Energy Resources) database, which provides satellite-derived and reanalysis-based climate records at a spatial resolution of $0.5° \times 0.625°$. The dataset was requested for the period January 1, 2020 through April 12, 2026 at the field coordinates (38.298°N, 48.847°E). The following variables were extracted: mean air temperature at 2~m (T2M, °C), maximum air temperature at 2~m (T2M\_MAX, °C), corrected daily precipitation (PRECTOTCORR, mm/day), all-sky surface shortwave downward irradiance (ALLSKY\_SFC\_SW\_DWN, MJ/m$^2$/day), root zone soil wetness (GWETROOT, dimensionless, 0--1), and surface soil wetness (GWETTOP, dimensionless, 0--1). The last two variables were used for validation purposes (Section~\ref{sec:validation}) rather than as model inputs.

Two days within the six-year record contained missing values (coded as $-999$ in the NASA POWER format). For temperature, precipitation, and soil wetness, the missing entries were replaced with the respective column means. For solar radiation, missing values were imputed using monthly averages to preserve seasonal patterns, since radiation exhibits strong seasonal variation.


\subsection{Reference Evapotranspiration}
\label{sec:et0}

Reference evapotranspiration (ET$_0$) is the primary driver of crop water demand. Since the NASA POWER dataset provides temperature and radiation but not the full suite of variables required for the Penman--Monteith equation (notably wind speed and relative humidity), ET$_0$ is estimated using the Hargreaves equation (FAO Irrigation and Drainage Paper No.~56, Equation~52):
\begin{equation}
    \text{ET}_0 = 0.0023 \cdot (T_{\text{mean}} + 17.8) \cdot (T_{\max} - T_{\min})^{0.5} \cdot R_s
    \label{eq:et0}
\end{equation}
where $T_{\text{mean}}$ is the mean daily temperature, $T_{\max}$ and $T_{\min}$ are the daily maximum and estimated minimum temperatures, and $R_s$ is the incoming solar radiation. The minimum temperature is estimated as $T_{\min} = 2T_{\text{mean}} - T_{\max}$, following standard practice when only mean and maximum temperatures are available. The Hargreaves equation is widely used in data-sparse regions and has been validated extensively for semi-arid and Mediterranean climates.


\subsection{Scenario Selection}
\label{sec:scenarios}

Two contrasting years were selected from the six-year record to represent the range of climatic conditions that the irrigation controllers must handle:

\paragraph{Dry scenario (2021).} The 120-day crop season beginning April~1, 2021 received only 90~mm of total rainfall, with nearly all precipitation concentrated in the first 10 days. The remaining 110 days were essentially rain-free. Mean temperature was 19.9°C with moderate variability, and mean ET$_0$ was 7.1~mm/day. This scenario represents an extreme drought year in which the irrigation controller bears almost full responsibility for maintaining crop water supply.

\paragraph{Wet scenario (2022).} The 120-day season beginning April~1, 2022 received 474~mm of rainfall, more than five times the dry scenario. However, this rainfall was heavily concentrated in a single extreme event during the first week (approximately 370~mm in 3 days), followed by a relatively dry remainder of the season. Mean temperature was 17.5°C and mean ET$_0$ was 5.5~mm/day. This scenario tests the controller's ability to avoid over-irrigation following a major rainfall event and to manage waterlogging risk in low-elevation agents.

\begin{figure}[htbp]
    \centering
    \includegraphics[width=0.85\textwidth]{figures/Figure_4.png}
    \caption{Comparison of the dry (2021, orange) and wet (2022, blue) climate scenarios over the 120-day crop season. Top left: daily rainfall---note the extreme precipitation event in the wet scenario. Top right: mean temperature. Bottom left: reference evapotranspiration. Bottom right: solar radiation.}
    \label{fig:climate}
\end{figure}

The selection of these two scenarios (Figure~\ref{fig:climate}) provides a demanding test of controller robustness. The dry scenario evaluates whether the controller can maximize biomass under severe water scarcity, while the wet scenario evaluates whether it can avoid wasting budget when natural rainfall is abundant and manage the waterlogging risk created by extreme precipitation events. Both scenarios use identical soil parameters and the same 130-agent field topology, isolating the climate as the sole source of variation.


% =============================================================================
\section{Soil and Crop Parameters}
\label{sec:soil_crop}
% =============================================================================

The soil at the study site is classified as silty loam, based on regional soil surveys and FAO soil classification maps for Gilan Province. The soil hydraulic properties used in this study follow FAO Irrigation and Drainage Paper No.~56 recommendations for this soil type. The root zone depth is set to $\theta_5 = 400$~mm, which is appropriate for irrigated lowland rice in the region.

\begin{table}[htbp]
\centering
\caption{Soil hydraulic parameters for silty loam (common to all crops).}
\label{tab:soil}
\begin{tabular}{clc}
\toprule
\textbf{Parameter} & \textbf{Description} & \textbf{Value} \\
\midrule
$\theta_1$ & Water uptake coefficient & 0.096 \\
$\theta_2$ & Wilting point (volumetric fraction) & 0.15 \\
$\theta_3$ & Initial abstraction threshold (mm) & 10.0 \\
$\theta_4$ & Drainage coefficient & 0.05 \\
$\theta_5$ & Root zone depth (mm) & 400 \\
$\theta_6$ & Field capacity (volumetric fraction) & 0.35 \\
\bottomrule
\end{tabular}
\end{table}

The resulting soil water thresholds that govern model behavior are:
\begin{itemize}
    \item Field capacity: $\text{FC} = \theta_6 \times \theta_5 = 0.35 \times 400 = 140$~mm
    \item Wilting point: $\text{WP} = \theta_2 \times \theta_5 = 0.15 \times 400 = 60$~mm
    \item Available water: $\text{AW} = \text{FC} - \text{WP} = 80$~mm
    \item Stress threshold (RAW): $\text{FC} - p \cdot \text{AW} = 140 - 0.20 \times 80 = 124$~mm
\end{itemize}
where $p = 0.20$ is the FAO depletion fraction for rice. The narrow band between the stress threshold (124~mm) and field capacity (140~mm) reflects rice's extreme sensitivity to soil water depletion: stress begins after only 20\% of the available water has been consumed. This places severe demands on the irrigation controller, which must maintain soil moisture within a 16~mm operating window to avoid triggering drought stress.

\begin{table}[htbp]
\centering
\caption{Crop-specific parameters for rice (primary crop).}
\label{tab:rice}
\begin{tabular}{clc}
\toprule
\textbf{Parameter} & \textbf{Description} & \textbf{Value} \\
\midrule
$\theta_7$    & Base temperature for phenology (°C)                     & 10.0 \\
$\theta_9$    & Heat stress onset temperature (°C)                      & 35.0 \\
$\theta_{10}$ & Extreme heat threshold (°C)                             & 42.0 \\
$\theta_{11}$ & Heat stress effect on maturity                          & 0.003 \\
$\theta_{12}$ & Drought stress effect on maturity                       & 0.003 \\
$\theta_{13}$ & Radiation use efficiency (g/MJ)                         & 1.20 \\
$\theta_{14}$ & Drought sensitivity on RUE                              & 0.80 \\
$\theta_{18}$ & Cumulative temperature to maturity (°C$\cdot$day)       & 1800 \\
$\theta_{19}$ & Maximum radiation interception fraction                  & 0.95 \\
$\theta_{20}$ & Thermal time for 50\% interception (°C$\cdot$day)       & 600 \\
$K_c$         & Season-average crop coefficient (FAO)                   & 1.15 \\
$p$           & Depletion fraction before stress (FAO)                  & 0.20 \\
\bottomrule
\end{tabular}
\end{table}

The selection of rice as the primary crop (Table~\ref{tab:rice}) was motivated by the initial finding that wheat ($\theta_{14} = 0.5$, $p = 0.55$) showed negligible biomass response to irrigation in the corrected ABM. With a deep stress threshold at only 96~mm and moderate drought sensitivity, wheat could sustain near-optimal growth from rainfall alone in both the dry and wet scenarios, rendering the controller comparison uninformative. Rice, with its narrow 16~mm operating window and severe drought penalty, provides the maximum contrast between irrigated and unirrigated outcomes and thus the most demanding test of controller performance.


% =============================================================================
\section{Water Budget Formulation}
\label{sec:water_budget}
% =============================================================================

A central feature of this thesis is the enforcement of a finite seasonal water budget, reflecting the reality of water scarcity in Iranian agriculture~\cite{nouri2023,mesgaran2018}. The budget computation follows the FAO crop water requirement method, which provides an agronomically grounded estimate of how much irrigation water the crop actually needs, rather than the arbitrary percentage-of-field-capacity approach used in some earlier studies.

The full seasonal irrigation requirement per agent is computed as the difference between crop water demand and natural rainfall supply:
\begin{equation}
    W_{\text{full}} = \max\!\bigl(\overline{\text{ET}_0} \cdot K_c \cdot n_{\text{days}} - R_{\text{total}},\; 0\bigr)
    \label{eq:wfull}
\end{equation}
where $\overline{\text{ET}_0}$ is the mean daily reference evapotranspiration over the season (mm/day), $K_c = 1.15$ is the FAO season-average crop coefficient for rice, $n_{\text{days}} = 120$ is the season length, and $R_{\text{total}}$ is the total seasonal rainfall (mm). The water budget per agent is then obtained by applying a scarcity factor $S \in (0, 1]$:
\begin{equation}
    W_{\text{agent}} = S \cdot W_{\text{full}}
    \label{eq:wagent}
\end{equation}
\begin{equation}
    W_{\text{total}} = W_{\text{agent}} \cdot N
    \label{eq:wtotal}
\end{equation}

The scarcity factor $S$ represents the fraction of the full crop water requirement that is actually available for irrigation, reflecting regional water allocation constraints. For the primary experiments in this thesis, $S = 0.70$ (70\% of full requirement) was adopted, representing moderate Iranian scarcity as documented in the literature on Iranian water allocation policy~\cite{afrasiabikia2026}. This choice means the controller must produce near-optimal biomass using only 70\% of the water the crop would ideally require, forcing genuine optimization rather than trivial saturation.

\begin{table}[htbp]
\centering
\caption{Water budget calculations for the two climate scenarios ($S = 0.70$).}
\label{tab:budget}
\begin{tabular}{lcc}
\toprule
\textbf{Variable} & \textbf{Dry (2021)} & \textbf{Wet (2022)} \\
\midrule
Mean ET$_0$ (mm/day)                  & 7.1   & 5.5   \\
Seasonal ET demand (mm/agent)         & 980   & 759   \\
Seasonal rainfall (mm)                & 90    & 474   \\
Full irrigation need (mm/agent)       & 890   & 285   \\
Budget per agent at $S=0.70$ (mm)     & 623   & 200   \\
Total field budget (mm)               & 80,990 & 26,000 \\
\bottomrule
\end{tabular}
\end{table}

The dry scenario requires more than three times the irrigation budget of the wet scenario (Table~\ref{tab:budget}), reflecting the dramatic difference in natural water supply. In both cases, the controller must allocate its budget efficiently across 130 agents over 120 days, accounting for the uneven spatial distribution of water caused by the field's topography.


% =============================================================================
\section{ABM Validation and Satellite Cross-Check}
\label{sec:validation}
% =============================================================================

Model validation was performed in two stages. First, the corrected ABM was run under a fixed open-loop irrigation schedule to verify that the soil water and biomass dynamics exhibit physically plausible behavior. Second, the simulated soil moisture trajectory was cross-checked against satellite-derived soil wetness observations from the NASA POWER dataset.


\subsection{Fixed-Schedule Validation}
\label{sec:fixed_validation}

The ABM was initialized with soil water at 60\% of field capacity ($x_1 = 84$~mm) across all agents, with all other states set to zero. A fixed irrigation schedule was applied under the dry (2021) climate scenario, with irrigation events occurring every 4 days. The per-event irrigation amount was computed by distributing the total budget (at $S = 0.50$ for this validation run) uniformly across the 30 irrigation events over the 120-day season.

\begin{figure}[htbp]
    \centering
    \includegraphics[width=0.85\textwidth]{figures/abm_validation.png}
    \caption{ABM validation under fixed-schedule irrigation (dry scenario, rice). Top: mean soil water across 130 agents with field capacity (140~mm), wilting point (60~mm), and stress threshold (124~mm) marked. Bottom: mean biomass accumulation showing the expected sigmoid (S-curve) growth pattern.}
    \label{fig:abm_validation}
\end{figure}

Figure~\ref{fig:abm_validation} demonstrates that the corrected ABM produces physically consistent behavior. The soil water state fluctuates between the stress threshold and field capacity, with periodic sawtooth patterns corresponding to irrigation events followed by transpiration-driven depletion. The soil water never drops below the wilting point, confirming that the irrigation schedule is adequate. The biomass curve exhibits the expected sigmoid shape, with a slow initial phase (days 0--40) corresponding to low radiation interception, followed by rapid exponential growth (days 40--90) as canopy cover expands, and a deceleration toward maturity (days 90--120). This S-curve is characteristic of the logistic growth function (Equation~\ref{eq:growth}) and qualitatively matches the biomass trajectories reported by Lopez-Jimenez et al.~\cite{Lopezjimenez2024}.

Several specific features confirm the physical correctness of the corrected model. First, soil water responds immediately to irrigation events (sharp upward jumps every 4 days) and depletes smoothly between events due to transpiration and drainage. Second, the soil water is bounded below field capacity and above the wilting point throughout the season, confirming that the mass balance equation and the drainage mechanism operate correctly. Third, the biomass curve begins its rapid growth phase around day~40, coinciding with the thermal time reaching the 50\% interception threshold ($\theta_{20} = 600$~°C$\cdot$day at approximately 10--15~°C/day accumulation). Fourth, the growth deceleration in the final 30 days reflects the senescence branch of the growth function as cumulative stress $x_3$ and thermal time $x_2$ approach the maturity threshold $\theta_{18}$.


\subsection{MPC Baseline Verification}
\label{sec:mpc_verification}

To further validate the simulation environment, the MPC controller was run under the dry scenario with prediction horizon $H_p = 8$ and water cost weight $\lambda = 0.05$. This serves as a verification that the complete control loop (ABM + optimizer + budget enforcement) operates correctly before the controller comparison experiments of Chapter~5.

\begin{figure}[htbp]
    \centering
    \includegraphics[width=0.85\textwidth]{figures/mpc_dry_rice_Hp8.png}
    \caption{MPC results for rice under the dry scenario (2021) with $H_p = 8$. Top: soil water trajectory maintained near the stress threshold. Middle: biomass growth with MPC (green) versus no-irrigation baseline (dashed gray), showing approximately 140\% improvement. Bottom: daily irrigation schedule, with the full budget exhausted by day~95.}
    \label{fig:mpc_validation}
\end{figure}

Figure~\ref{fig:mpc_validation} demonstrates several important properties of the MPC-controlled system. The soil water is maintained in the optimal range between the stress threshold and field capacity for the first 90 days, after which the budget is exhausted and soil water declines. The MPC front-loads irrigation during the critical vegetative and early reproductive phases, which aligns with agronomic best practice for rice. The biomass trajectory shows a clear and substantial improvement over the no-irrigation baseline, confirming that the drought stress correction (Section~\ref{sec:h3_correction}) successfully makes irrigation consequential for yield outcomes. The irrigation schedule reveals that the optimizer distributes water across nearly every day of the season rather than applying large intermittent doses, reflecting the optimizer's preference for maintaining soil moisture near the stress threshold to minimize drought penalty.


\subsection{Satellite Soil Wetness Cross-Validation}
\label{sec:satellite}

As an independent validation check, the simulated soil moisture was compared against the NASA POWER root zone soil wetness (GWETROOT) product. GWETROOT is a dimensionless index (0--1) derived from the MERRA-2 reanalysis, representing the fraction of saturated soil moisture in the root zone. While this product has coarser spatial resolution than the agent grid and reflects large-scale rather than field-scale conditions, it provides a useful qualitative check on the seasonal pattern of soil moisture evolution.

The comparison confirmed that the simulated soil moisture trajectory follows the same seasonal pattern as the satellite observations: both show high initial moisture in early April (consistent with spring snowmelt and early rainfall), a sustained decline through the dry summer months, and partial recovery following late-season rainfall in the wet scenario. The magnitude of seasonal depletion in the simulation is consistent with the satellite data, supporting the validity of the transpiration, drainage, and runoff parameterizations. A detailed quantitative comparison would require sub-field-scale soil moisture observations, which are not available from satellite products at this resolution; however, the qualitative agreement provides confidence that the ABM captures the essential features of the soil water balance at this site.


% =============================================================================
\section{Chapter Summary}
\label{sec:ch3_summary}
% =============================================================================

This chapter has presented the complete system description and modeling framework for the irrigation control study. The key elements can be summarized as follows:

\begin{enumerate}
    \item \textbf{Study site.} A 6-hectare field in Gilan Province, Iran, discretized into 130 crop--soil agents on a $10 \times 13$ grid with a 111~m elevation range and a bowl-shaped topology.
    
    \item \textbf{Agent-based model.} Four coupled difference equations per agent capture soil water dynamics, thermal accumulation, crop maturation, and biomass growth, with lateral water exchange governed by a topography-derived directed graph.
    
    \item \textbf{Original corrections.} Two modifications to the baseline model of Lopez-Jimenez et al.~\cite{Lopezjimenez2024}: (a) a mass-conserving water exchange formulation that subtracts outflow from the sender, and (b) the inclusion of drought stress $h_3$ as a direct multiplier in the biomass equation.
    
    \item \textbf{Climate scenarios.} NASA POWER data provide a dry scenario (2021, 90~mm rainfall) and a wet scenario (2022, 474~mm rainfall) spanning the range of conditions the controllers must handle.
    
    \item \textbf{Crop selection.} Rice ($p = 0.20$, $\theta_{14} = 0.8$, $K_c = 1.15$) was selected as the primary crop due to its extreme water sensitivity, which provides the most demanding test case for controller comparison.
    
    \item \textbf{Water budget.} The FAO $K_c$ method with a 70\% scarcity factor provides an agronomically grounded seasonal budget reflecting Iranian water allocation constraints.
    
    \item \textbf{Validation.} Fixed-schedule simulations and satellite soil wetness cross-checks confirm that the corrected ABM produces physically consistent behavior.
\end{enumerate}

With the simulation environment fully specified and validated, the next chapter proceeds to the design of the three controllers that will be compared: the threshold baseline, the MPC formulation, and the SAC-RL agent.